\originalchapter{进阶电磁理论}
\originalsection{色散与因果响应}
8.311 的课历把介质色散放在波传播之前, 因为实际物质响应不是瞬时的. 对电子位移 $x$ 取受迫阻尼振子模型 $\ddot x+\Gamma\dot x+\omega_0^2x=-(e/m)E(t)$. 对 $e^{-i\omega t}$ 约定, 单位体积 $N$ 个电子给
\[
\chi_e(\omega)=\frac{Ne^2}{\eps m}\frac1{\omega_0^2-\omega^2-i\Gamma\omega}.
\]
极化为 $P=-Nex$, 电子电荷的两个负号使静态极化率为正. 阻尼 $\Gamma>0$ 且正频率 $\omega>0$ 时 $\operatorname{Im}\chi_e>0$ 对应吸收, 平均吸收功率密度为 $\eps\omega\operatorname{Im}\chi_e|E_0|^2/2$; 负频率的虚部反号; 无阻尼极限的共振须用分布或极限处理, 不能把极点处表达式当有限响应.

更一般地写 $P(t)=\eps\int_0^\infty\chi(\tau)E(t-\tau)\dd\tau$, 下限0表达因果性. 假定 $\chi(\omega)\to0$ 的高频行为足以使无减除积分收敛; 若趋向常数 $\chi_\infty$, 应对 $\chi-\chi_\infty$ 写关系. 在因果核适当衰减的条件下, 其傅里叶响应在上半复频平面解析. Kramers--Kronig 关系把实部与虚部联系为
\[
\operatorname{Re}\chi(\omega)=\frac1\pi\PV\int_{-\infty}^\infty
\frac{\operatorname{Im}\chi(\omega')}{\omega'-\omega}\dd\omega'.
\]
这是调用复分析的色散关系, 本稿不把它当成已从零证明的本科定理. 正负频率的关系由实响应 $\chi(-\omega)=\chi(\omega)^*$ 给出. 介电常数与吸收不能随意独立选择, 否则可能违反因果性.

\originalsection{负折射率与切伦科夫辐射}
对有效线性、各向同性、均匀介质在单频透明频带中的平面波, $\vct H=(\vct k\times\ee)/(\mu\omega)$, 所以平均能流方向与 $\vct k/\mu$ 相同. 当一定频带内 $\varepsilon<0$ 且 $\mu<0$ 时, 能流与相位波矢反向. 这不表示能量逆时间传播, 还须用物质色散和正储能决定物理分支. 8.07 第12、13讲讨论的超材料是这种进阶例子, 不能将负常数 $\varepsilon,\mu$ 无频带限制地外推.

匀速粒子在介质中的场含 $\omega=\vct k\cdot\vct v$ 的相位匹配条件. 若普通正折射传播分支的相速度为 $c/n(\omega)$, 则各频率的切伦科夫条件为 $v>c/n(\omega)$、$\cos\theta_k(\omega)=c/[n(\omega)v]$. 色散使锥角随频率变化. 在负折射介质中必须区分波矢的相位锥与能流的辐射锥; 用正折射公式只换 $n$ 的负号而不说明角的方向会造成误导. 本稿只建立相位与能流关系, 不复用研究报告中的第三方实验图.

\originalsection{传输线与场振子}
假设填充介质为静止、线性、均匀、各向同性. 同轴线内外半径 $a,b$, 单位长度电容 $C'=2\pi\varepsilon/\ln(b/a)$, 单位长度电感 $L'=\mu\ln(b/a)/(2\pi)$. 前者来自径向电场, 后者来自 $B_\phi=\mu I/(2\pi s)$ 的磁能 $U_B/\ell=L'I^2/2$. 因而传播速度 $1/\sqrt{L'C'}=1/\sqrt{\mu\varepsilon}$, 阻抗 $\sqrt{L'/C'}=\sqrt{\mu/\varepsilon}\ln(b/a)/(2\pi)$. 这是把分布场和电路变量对应起来, 不是把整条线压缩成单个电容.

取无源真空理想导体腔, 无静电背景, 并选择相容的辐射规范 $V=0$. 理想腔中的矢量势可按正交横向模态展开 $\vct A=\sum_\alpha q_\alpha(t)\vct f_\alpha(\rr)$, 库仑规范要求 $\nabla\cdot\vct f_\alpha=0$. 模态满足相应导体边界及 $-\nabla^2\vct f_\alpha=(\omega_\alpha/c)^2\vct f_\alpha$. 选择 $\eps\int\vct f_\alpha\cdot\vct f_\beta\dd^3x=\delta_{\alpha\beta}$ 归一化后, 真空场拉格朗日量成为 $L=\frac12\sum_\alpha(\dot q_\alpha^2-\omega_\alpha^2q_\alpha^2)$, 哈密顿量成为独立谐振子能量之和. 8.311 的 \textit{Field Oscillators} 补充笔记由此进入量子化前的经典基础; 本稿到经典模态为止.

\originalsection{协变与证明范围}
沿 $x$ 方向的洛伦兹变换为 $ct'=\gamma(ct-\beta x)$、$x'=\gamma(x-\beta ct)$、$y'=y$、$z'=z$. 平行场分量不变, 垂直分量满足 $\ee'_\perp=\gamma(\ee+\vct v\times\bb)_\perp$、$\bb'_\perp=\gamma(\bb-c^{-2}\vct v\times\ee)_\perp$. 这里新系相对旧系以 $+\vct v$ 运动, 换回旧系要反转速度. 不变量为 $E^2-c^2B^2$ 与 $\ee\cdot\bb$.

源与势可写四维量, 麦克斯韦方程可写反对称场张量方程. 本章只给与本科题相接的变换、色散和模态入口; Noether 定理、完整协变作用量、辐射反作用及同步辐射的全套推导尚不在这一本科重编的证明范围. 它们有对应的 8.311 原始作业与官方解答, 见后面的进阶资料定位. 不把研究生课程课历中的题目清单写成已经完成的理论章节.
